\documentclass[12pt]{article}
\usepackage[utf8]{inputenc}
\usepackage[T1]{fontenc}
\usepackage[margin=1in]{geometry}
\usepackage{setspace}
\usepackage{xcolor}
\usepackage[colorlinks=true, urlcolor=blue!60!black, linkcolor=blue!60!black]{hyperref}
\setlength{\parindent}{0em}
\setlength{\parskip}{0.6em}
\renewcommand{\baselinestretch}{1.0}
\begin{document}
\begin{center}
{\LARGE \textit{Rucho v. Common Cause}, 588 U.S. 684 (2019)\par}
\vspace{0.5em}
{\large Case Brief --- Two-Page Edition\par}
\end{center}

\section{Memory Jogger}
Federal courts cannot hear partisan gerrymandering claims: no manageable, neutral standard tells how much partisanship is too much (\href{https://www.courtlistener.com/opinion/4633469/rucho-v-common-cause/}{CourtListener}).

\section{Detailed Case Facts}
North Carolina Republicans drew a 2016 congressional map with ``Partisan Advantage'' as an express criterion, aiming for 10 of 13 seats. Maryland Democrats redrew the Sixth District in 2011, moving some 700,000 voters to flip it and secure 7 of 8 seats.

\section{Procedural History}
Three-judge courts struck down both maps---North Carolina's under equal protection, the First Amendment, and Article I; Maryland's under the First Amendment. The states appealed directly.

\section{Judicial Votes}
Vacated 5--4. Roberts wrote for Thomas, Alito, Gorsuch, and Kavanaugh; Kagan, joined by Ginsburg, Breyer, and Sotomayor, dissented.

\section{Holding}
Partisan gerrymandering claims under any constitutional provision are nonjusticiable political questions; the cases must be dismissed for lack of jurisdiction. Remarks about state and congressional reform are dicta.

\section{Analysis of Opinions}
\textbf{Majority:} Some partisanship is permissible, so the question is ``How much is too much?'' Claims rest on proportional representation, which the Constitution does not require; predicting durability asks judges to ``beat the point spread''; maps restrict no speech. Unlike population equality and racial classifications, partisanship offers no legal standard. \textbf{Kagan:} gerrymanders let politicians choose voters; lower courts had a workable test using simulated maps to find ``extreme outliers''---``This much is too much.''

\section{Examples: Future Applications}
\textbf{Same-Side:} a software-drawn state senate map locking in a supermajority is unreviewable in federal court. \textbf{Opposite-Side:} state courts may enforce state constitutional limits (\textit{Moore v. Harper}). \textbf{Fence-Sitter:} moving Black Democratic voters where race and party correlate.

\section{Critique}
John McGinnis sees ``silent originalism'' that needed no political question label. Election-law scholars call the decision an abdication where politics cannot self-correct, one that greenlights technological entrenchment.

\section{Key Quotations}
\begin{quote}
``Excessive partisanship in districting leads to results that reasonably seem unjust. But the fact that such gerrymandering is `incompatible with democratic principles,' \textit{Arizona State Legislature}, 576 U.\,S., at 791 (brackets omitted), does not mean that the solution lies with the Federal Judiciary. We conclude that partisan gerrymandering claims present political questions beyond the reach of the federal courts.''
\end{quote}
\noindent --- Chief Justice Roberts, for the Court, 588 U.S. at 718.

\end{document}
